\documentclass[11pt]{amsart}
\usepackage{amsmath,amssymb,amsthm,mathtools}
\usepackage[margin=1in]{geometry}
\usepackage{enumerate}
\usepackage{hyperref}

\newtheorem{theorem}{Theorem}[section]
\newtheorem{proposition}[theorem]{Proposition}
\newtheorem{lemma}[theorem]{Lemma}
\newtheorem{corollary}[theorem]{Corollary}
\theoremstyle{definition}
\newtheorem{definition}[theorem]{Definition}
\newtheorem{principle}[theorem]{Principle}
\theoremstyle{remark}
\newtheorem{remark}[theorem]{Remark}
\newtheorem{observation}[theorem]{Observation}

\newcommand{\RH}{\mathrm{RH}}
\newcommand{\bbR}{\mathbb{R}}
\newcommand{\bbC}{\mathbb{C}}
\newcommand{\bbN}{\mathbb{N}}
\newcommand{\bbZ}{\mathbb{Z}}
\newcommand{\Koff}{\mathcal{K}_{\mathrm{off}}}
\newcommand{\negind}{\operatorname{neg.ind}}
\newcommand{\Wtwo}{\mathcal{W}_2}
\newcommand{\dnu}{d\nu}
\newcommand{\Tl}{T_\lambda}
\DeclareMathOperator{\Real}{Re}
\DeclareMathOperator{\sig}{sig}
\DeclareMathOperator{\Spec}{Spec}
\newcommand{\La}{\Lambda}

\title[One Wall in Many Coordinates]{One wall in many coordinates:\\
closure of the propagation trichotomy\\
and the master quantifier for the Riemann Hypothesis}

\author{David Alejandro Trejo Pizzo}
\thanks{Independent Researcher. \texttt{dtrejopizzo@gmail.com}}


\begin{document}
\nocite{bgTitchmarsh,bgEdwards,bgConrey03,bgBombieri00,bgConnes99,bgdeBranges,bgMont73,bgBK}

\begin{abstract}
The companion papers \cite{P41,P42} reduced the search for a proof of the
Riemann Hypothesis (RH), within the Connes--Consani--Moscovici framework, to
a trichotomy of argument-forms surviving a structural no-go: two-dimensional
propagation (Form~A), index amplification (Form~B), and an effective
self-approximation criterion (Form~C). All three are now closed. We collect
the three closure theorems and make the paper's central claim: they are not
three independent failures but one obstruction expressed in three coordinate
systems. In each case the unconditionally controlled object is an
\emph{averaged}, \emph{generic}, or \emph{interior} quantity, and RH is
equivalent to an \emph{individual}, \emph{uniform}, or \emph{boundary}
refinement of it; the wall is always the inversion of a single quantifier,
from ``on average / for almost all / in the interior'' to ``for the specific
/ uniformly / at the boundary.'' We show this quantifier inversion is the
same object as the parity barrier that closed the $\omega$-class arc, the
arithmetic-propagation wall MW-2, the Hadamard barrier MW-7, and the
individual-versus-generic orbit obstruction of Bagchi's criterion. The
sharpest new formulation, isolated in the amplification analysis, is a
distinction between two kinds of autonomy of a quadratic functional: its
\emph{value} can be computed from the primes alone, but its \emph{inertia}
(the partition of its negative index into blocks) cannot be determined
without the positions of the zeros. We state the master quantifier as a
synthesis principle, exhibit each closure as an instance, and present the
unconditional ledger the programme leaves behind---among them the
strict single-scale trace criterion, the tensor index identity
$\negind(Q_1\otimes Q_2)=p_1q_2+q_1p_2$, an unconditional sub-resolution
positivity theorem, and a quadruple-density target weaker than RH. We close
with an honest account of what a proof must now do: cross the mean-to-uniform
barrier by a mechanism no current argument-form provides.
\end{abstract}

\maketitle
\tableofcontents
\newpage

\section{Introduction}

\subsection{The end of a programme, and what it leaves}
A long programme \cite{P40,P41,P42} set out to prove RH and did not. What it
produced instead is a map: a sequence of attempts, each documented and each
either refuted or reduced to a named wall, culminating in the structural
no-go of \cite{P41}. That paper argued that RH always appears as the
boundary value of a one-parameter family whose interior is an unconditional
theorem, that no one-dimensional propagation principle reaches the boundary,
and that exactly three argument-forms survive: two-dimensional propagation
(Form~A), index amplification (Form~B), and an effective version of Bagchi's
self-approximation criterion (Form~C). The companion paper \cite{P42}
pursued all three. Forms~A and~C closed there; Form~B survived a first
decisive experiment and was reduced to a single missing lemma, the
sieve-to-index bridge. That bridge has since closed \cite[]{P40}.

This paper does two things. First (\S\ref{sec:trichotomy}), it records the
three closures as a single ledger of theorems. Second, and principally
(\S\ref{sec:master}--\S\ref{sec:autonomy}), it argues that the three are one
wall in three coordinate systems, isolates that wall as a single quantifier
inversion, and gives its sharpest formulation. The remaining sections present
the unconditional results the programme leaves behind
(\S\ref{sec:ledger}) and an honest account of the boundary of current
methods (\S\ref{sec:boundary}).

\subsection{What is claimed}
We claim no proof of RH and no conditional result following from a weaker
hypothesis than is already known. The central claim is structural and is
stated as a synthesis principle (Principle~\ref{prin:master}), not a theorem:
that the obstruction to RH, across every arc of the programme, is the
inversion of one quantifier. Several \emph{instances} of the principle are
theorems (the closure theorems of \S\ref{sec:trichotomy}); the principle
itself is the assertion that they are the same theorem. The one genuinely new
conceptual contribution is the value/inertia distinction of
\S\ref{sec:autonomy}, which we believe is the most precise form yet of the
Hadamard barrier.

\subsection{Framework}
We use the notation of \cite{P39,P40,P41,P42}: $\dnu\ge0$ is the
unconditional discrepancy measure with $\RH\iff\dnu=0$; the bridge theorem
\cite{P35} gives $\kappa(Q)=\negind(H_C)=2m$ on the off-critical subspace
$\Koff$, with $m$ the number of off-critical quadruples and $\RH\iff m=0$;
$\Wtwo$ is the two-variable Weil functional of \cite[]{P40}.

\section{The trichotomy, closed}\label{sec:trichotomy}

We restate the three closures as theorems, citing their proofs. Each
identifies the precise point at which the form fails.

\begin{theorem}[Form~C: the self-blocking amplifier {\cite[]{P40}, \cite[\S2]{P42}}]
\label{thm:C}
The self-approximation density $\delta(d,\varepsilon)$ does not degenerate as
$d\to0$; the guaranteed scale degenerates linearly,
$T_0\asymp C(\varepsilon)/|d|$. The Rouch\'e amplification that would convert
one off-critical zero into a positive-density family is bounded, under
$\neg\RH$, by an anchored density carrying the \emph{same} exponent
$\kappa(\sigma)$ as the unconditional zero-density ceiling
$N(\sigma,T)=O(T^{\kappa(\sigma)})$ that would contradict it. The density
needed to force RH and the ceiling that caps it move together under any
improvement; mobile $d=d(T)$ does not separate them.
\end{theorem}

\begin{theorem}[Form~A: degeneration {\cite[]{P40}, \cite[\S3]{P42}}]
\label{thm:A}
Let $\mathcal{T}(t,x)$ be the joint heat-time/arithmetic discrepancy trace.
The only definition for which both edges---vanishing for $t\ge\Lambda$ and
reality of zeros at finite $x$---are unconditional theorems vanishes
identically on the interior of the quadrant, because the forward
De~Bruijn--Newman flow preserves reality of zeros \cite{deBruijn50}. Hence
$\partial_t\mathcal{T}\equiv\partial_x\mathcal{T}\equiv0$ and every coupling
inequality is vacuous; the corner $(0,\infty)=\RH$ remains a pure
discontinuity.
\end{theorem}

\begin{theorem}[Form~B: inertia is not autonomous {\cite[]{P40}}]
\label{thm:B}
The negative index amplifies exponentially and unconditionally,
$\kappa(\psi^{\otimes k})=\tfrac12(4m)^k$ on $\Koff$
\cite[]{P40}, and the tensor square is realised by the two-variable
Weil functional $\Wtwo$ on a thickened multiplicative diagonal, whose value
is computable from the primes alone \cite[]{P40}. Nevertheless the
window indices of $\Wtwo$ are binary---$0$ on every window under RH, and
$\asymp m\log T$ without bound under $\neg\RH$---so any uniform finite bound
on the relative negative index already implies RH at $k=2$, before
amplification can act. The bound the bridge requires is an asymptotic
evaluation of the almost-prime functional with relative error $o(1/\log T)$
uniform over all windows, which implies RH directly.
\end{theorem}

The three forms fail at three apparently different places: a matching of
density exponents (C), an invariance of the heat flow (A), a binarity of
window indices (B). The rest of the paper argues they are the same place.

\section{The master quantifier}\label{sec:master}

\subsection{The common shape}
In each closure, separate the object that is unconditionally controlled from
the object equivalent to RH.

\begin{center}
\renewcommand{\arraystretch}{1.45}
\begin{tabular}{lll}
\textbf{Form} & \textbf{Unconditionally controlled} & \textbf{Equivalent to RH}\\
\hline
C & zero density \emph{on average}, $N(\sigma,T)=O(T^{\kappa})$ & anchored density in the \emph{individual} window $J_d$\\
A & the trace on the \emph{interior} / forward flow & the trace at the \emph{boundary} corner $(0,\infty)$\\
B & the functional value \emph{in mean} over windows & the index \emph{uniformly} over every window\\
\end{tabular}
\end{center}

In all three the controlled quantity is an average, a generic value, or an
interior value; the RH-equivalent quantity is the corresponding individual,
specific, or boundary refinement. The gap between the two columns is, in each
row, the inversion of one quantifier.

\begin{principle}[The master quantifier]\label{prin:master}
Across the programme, RH is equivalent to a statement obtained from an
unconditional theorem by inverting a single quantifier:
\[
\text{``for almost all / on average / in the interior''}
\;\longrightarrow\;
\text{``for the specific / uniformly / at the boundary.''}
\]
Every unconditional method certifies the averaged side; none certifies the
refined side, because the arithmetic object sits in the measure-zero
exceptional set of its own generic behaviour, in the favourable direction,
and genericity cannot detect a measure-zero event.
\end{principle}

We do not claim Principle~\ref{prin:master} as a theorem; it is the
synthesis the closures support. Its force is that it is not specific to the
trichotomy: it organises the entire programme.

\subsection{The principle organises the whole programme}
The same quantifier inversion is visible in every arc.

\emph{Arc~A (the parity barrier).} Sieve methods control the count of
integers with given congruence data \emph{on average over residue classes};
RH-strength statements need to \emph{distinguish} $\omega(n)\bmod2$ for the
individual sequence. Selberg's parity principle is exactly the statement that
the averaged data cannot see the individual sign \cite[Arc~A]{P40}.

\emph{Arc~B (the arithmetic-propagation wall MW-2).} The Euler product
controls $\zeta$ \emph{generically} in the half-plane $\Real(s)>1$; the zeros
are \emph{individual} global objects in the strip, owned by no prime. The
information does not propagate because it is averaged information about a
region, not pointwise information about a locus.

\emph{The CCM Hadamard barrier MW-7.} No property of $\zeta$ verifiable
without the positions of the zeros---i.e.\ no \emph{generic}, position-blind
property---implies $\dnu=0$, the \emph{individual} vanishing.

\emph{Bagchi's criterion.} The Bohr-torus flow is uniquely ergodic, so the
\emph{generic} orbit recurs; RH is the recurrence of the \emph{individual}
arithmetic orbit $\omega_0$, which lies in the measure-zero exceptional set
\cite[]{P40}.

\begin{observation}
The four obstructions MW-2, MW-7, the parity barrier, and the Bagchi
individual-orbit gap, previously catalogued as distinct walls, are four
instances of Principle~\ref{prin:master}. The trichotomy closures
(Theorems~\ref{thm:C}--\ref{thm:B}) are three more. The programme's seven
walls collapse to one.
\end{observation}

\subsection{Why the inversion is hard in exactly this way}
The inversion is not a technical gap to be narrowed by sharper estimates. An
averaged statement and its individual refinement differ by the behaviour on a
null set, and \emph{the favourable direction of RH is precisely a null-set
event}: that the off-critical count is not merely $o(\text{generic})$ but
exactly zero. No method that proceeds by majorisation, equidistribution, or
density---all of which see only the averaged side---can certify an exact-zero
individual event. This is why every unconditional bound in the programme has
``the wrong sign'' (MW-4): it bounds the average, and the average is
compatible with both RH and its negation. The quantifier, not the size of
the constant, is the obstruction.

\section{Autonomy of value versus autonomy of inertia}\label{sec:autonomy}

The amplification analysis produced the sharpest formulation of the wall, and
we isolate it because it is the paper's one new conceptual tool.

\begin{definition}
Let $\Wtwo$ be the two-variable Weil functional and $Q_2$ its associated
quadratic form on the thickened-diagonal test space. We say the \emph{value}
of $Q_2$ is autonomous if $Q_2(h,h)$ is computable from arithmetic data
(primes, pole, gamma factor) without the positions of the zeros of $\zeta$.
We say the \emph{inertia} of $Q_2$ is autonomous if the partition of its
negative index into the relative block $\negind(Q|_{\Koff}^{\otimes2})$ and
the remainder is so computable.
\end{definition}

\begin{theorem}[Value autonomous, inertia not {\cite[--108]{P40}}]
\label{thm:autonomy}
The value of $Q_2$ is autonomous: the mixed zeros$\times$primes terms
reabsorb into arithmetic$\times$arithmetic terms, and the primes$\times$primes
block is the almost-prime functional
$\sum_N g(N)\sum_{d\mid N}\La(d)\La(N/d)\varphi(d^2/N)$, computable from the
primes alone. The inertia of $Q_2$ is \emph{not} autonomous: the negative
directions live, at second order, in the deficit between the almost-prime
kernel and its archimedean prediction---a pair-correlation quantity---and the
isolation of the off-critical block from the off$\times$on contamination
requires knowing which zeros are off the line. The global negative index is
$\infty$ under $\neg\RH$; only the relative block is finite, and the relative
block is not test-accessible.
\end{theorem}

\begin{remark}[The Hadamard barrier, sharpened]
Theorem~\ref{thm:autonomy} is the Hadamard barrier MW-7 in its most precise
form. Earlier statements asserted that no position-blind property implies
$\dnu=0$. The value/inertia distinction locates the obstruction one level
finer: one \emph{can} write down a position-blind functional whose value
encodes the discrepancy (the barrier does not forbid that), but one
\emph{cannot} read its index---the very invariant that detects $m$---without
the partition into blocks, and the partition is position-dependent. The
quantity is autonomous; the \emph{signature} is not. This is why the
amplification route, which needs the index and not the value, meets the wall
even though the value is in hand.
\end{remark}

\begin{observation}
Value/inertia is the master quantifier in operator-theoretic coordinates: the
value is the averaged (trace-like) datum; the inertia is the individual
(spectral-resolution) datum. Computing a trace is averaging; resolving a
signature is individuating. The same null-set obstruction reappears: the
off-critical contribution to the value is $o(1/\log T)$ per window---generic
zero---while its contribution to the inertia is a full negative direction---an
individual event the average cannot see.
\end{observation}

\section{The unconditional ledger}\label{sec:ledger}

A programme that does not reach its goal is judged by what it proves along the
way. We list the unconditional results, with sources, that survive
independently of RH.

\begin{enumerate}[\rm(L1)]
\item \emph{Non-negativity of the discrepancy:} $\dnu\ge0$ as a measure,
each off-critical quadruple contributing a factor $\ge1$
\cite[]{P40}.
\item \emph{Single-scale trace criterion:} the Abel kernel $W_\lambda$ is
strictly positive at every real point, whence $\RH\iff\Tl[\lambda_0]=0$ for a
single $\lambda_0$; the universal quantifier of the criterion is redundant
\cite[]{P40}, \cite[Thm~3.2]{P41}.
\item \emph{Tensor index identity:} for indefinite quadratic forms of
signatures $(p_i,q_i)$, $\negind(Q_1\otimes Q_2)=p_1q_2+q_1p_2$, strictly
superadditive; hence exponential amplification $\kappa(\psi^{\otimes k})=
\tfrac12(4m)^k$ on $\Koff$ \cite[]{P40}, \cite[Thm~4.2--4.3]{P42}.
\item \emph{Realisation of the tensor square:} the two-variable Weil
functional on the thickened diagonal realises
$(\Koff,Q|_{\Koff})^{\otimes2}$ with the predicted $(8,8)$ signature, the
thickening weight carrying the index \cite[]{P40}, \cite[Thm~5.4]{P42}.
\item \emph{Sub-resolution positivity:} $Q_2>0$ on test functions of fixed
support in high windows; the negative index is invisible below the resolution
scale $a\asymp\log T$ \cite[Thm~3.3]{P40}.
\item \emph{Closed-form kernel:} for $N(\lambda_0)=1$,
$B_{\lambda_0}(r)=\tfrac18\operatorname{sech}^{1/2}(r)
(105\operatorname{sech}^4 r-78\operatorname{sech}^2 r+5)$; its prime series
diverges like $(5\sqrt2/8)\log X$, so the CCM test function sits on the
boundary of the Weil class \cite[]{P40}, \cite[\S3]{P42}.
\item \emph{Eight equivalent reformulations} of RH as functionals of $\dnu$
\cite[]{P40}, and the bridge theorem $\kappa(Q)=\negind(H_C)$
\cite{P35}.
\item \emph{A quadruple-density target:} a bound of the form
$\kappa_W\le C\cdot n_W$ on window indices would give $m\le C/2$, an
unconditional bound on the number of off-critical quadruples weaker than RH
\cite[\S8.2]{P40}---behind the same wall, but a genuine partial
target.
\end{enumerate}

These are not consolation. (L2) and (L3) in particular are clean, quotable
facts; (L5) is an unconditional positivity with a precise resolution
threshold; (L8) is the shape of a theorem strictly weaker than RH that the
framework could plausibly reach.

\section{The boundary of current argument-forms}\label{sec:boundary}

\subsection{What the closures establish jointly}
\cite{P41} predicted that if Forms~A and~B also collapsed, the collapse would
itself be informative: it would be the first evidence that RH is unreachable
by any of the argument-forms currently known to mathematics. We have reached
that point, and the evidence is now specific rather than impressionistic.
Every form---propagation (1D in \cite{P41}, 2D in Form~A), amplification
(Form~B), recurrence (Form~C and Bagchi), positivity (MW-1, the whole of
Arc~B)---terminates at the same quantifier inversion. The walls are not many;
the wall is one, and we can now write it down.

\begin{principle}[The boundary]\label{prin:boundary}
A proof of RH must certify an individual, uniform, or boundary statement that
its averaged, generic, or interior shadow leaves undetermined---an exact-zero
event on a null set of the relevant generic behaviour. No argument-form
catalogued in this programme (propagation, amplification, recurrence,
positivity, sieve) does this, because each proceeds through the averaged
shadow. The missing idea must cross the mean-to-uniform barrier by a mechanism
outside this list.
\end{principle}

\subsection{Where such a mechanism might live}
We will not pretend to know the mechanism. We can say what property it must
have, sharpened by Theorem~\ref{thm:autonomy}: it must read the
\emph{inertia} of an arithmetic functional, not merely its value---it must
individuate a signature where only a trace is unconditionally available. In
the one historical case where RH-type statements were proved (function
fields), the mechanism that did this was cohomological: finite-dimensional
$\ell$-adic cohomology supplies an integer signature directly, bypassing the
averaged analytic data entirely. The programme's repeated arrival at the
arithmetic site, the scaling-site square, and the Riemann--Roch analogue
\cite[MW-5]{P40} is, in this light, not a coincidence but a recognition: a
genuinely cohomological invariant over $\Spec\bbZ$ would supply the inertia
non-generically, and that---not any analytic refinement---is what the master
quantifier demands. Whether such an invariant exists is the Connes--Consani
programme's open question, and it is outside the analytic methods this
programme exhausted.

\section{Assessment}\label{sec:assessment}

The programme set out to prove RH and did not; this paper records why, as
precisely as we are able. The structural trichotomy of \cite{P41} is closed
in all three branches (Theorems~\ref{thm:C}--\ref{thm:B}), and the three
closures, with the programme's earlier walls, are one obstruction: the
inversion of a single quantifier from averaged to individual
(Principle~\ref{prin:master}), in its sharpest form the distinction between
the autonomous value and the non-autonomous inertia of a quadratic functional
(Theorem~\ref{thm:autonomy}). The programme leaves an unconditional ledger
(\S\ref{sec:ledger}) and a clear statement of what a proof must do
(Principle~\ref{prin:boundary}).

We state the honest conclusion plainly. We have not narrowed the distance to
a proof of RH; the distance is, if anything, more visibly large than when we
began. What we have narrowed is the \emph{description} of the obstruction:
from a list of seven walls to a single quantifier inversion with a named
operator-theoretic form. A reader who hoped for progress toward the theorem
will find none here. A reader who wants to know exactly where the difficulty
lives, and why every analytic, spectral, and combinatorial route returns to
the same place, will find that this is now answerable in one sentence---and
that the answer points outside the methods we used, toward an invariant that
reads inertia rather than value. Mapping the wall is not crossing it. But a
wall whose single gate is known is a better thing to leave behind than a
fog.

\begin{thebibliography}{99}

\bibitem{CCM24} A.~Connes, C.~Consani, H.~Moscovici,
\emph{Spectral realisation of zeta zeros and the Riemann--Roch programme}
(2024).

\bibitem{deBruijn50} N.~G.~de~Bruijn, \emph{The roots of trigonometric
integrals}, Duke Math. J. \textbf{17} (1950), 197--226.

\bibitem{P35} D.~A.~Trejo Pizzo, \emph{Pontryagin spectral index of the Weil form and the Connes scaling operator: a bridge theorem for the Riemann Hypothesis}, preprint, 2026.

\bibitem{P39} D.~A.~Trejo Pizzo, \emph{A trace and measure criterion for the Riemann Hypothesis from the zeta spectral triple}, preprint, 2026.

\bibitem{P40} D.~A.~Trejo Pizzo, \emph{The localized Weil form and its Kre\u\i n--Pontryagin index: localization, obstruction, and rigidity}, preprint, 2026.

\bibitem{P41} D.~A.~Trejo Pizzo, \emph{The exhaustion of one-dimensional
propagation} (programme paper P41, 2026).

\bibitem{P42} D.~A.~Trejo Pizzo, \emph{Two closures and a survivor: the
amplification route to the Riemann Hypothesis through almost-prime sieve
bounds} (programme paper P42, 2026).

\bibitem{bgTitchmarsh} E.~C.~Titchmarsh, \emph{The Theory of the Riemann Zeta-Function}, 2nd ed., rev.\ D.~R.~Heath-Brown, Oxford Univ.\ Press, 1986.
\bibitem{bgEdwards} H.~M.~Edwards, \emph{Riemann's Zeta Function}, Academic Press, 1974.
\bibitem{bgConrey03} J.~B.~Conrey, \emph{The Riemann Hypothesis}, Notices Amer.\ Math.\ Soc.\ \textbf{50} (2003), 341--353.
\bibitem{bgBombieri00} E.~Bombieri, \emph{Problems of the Millennium: The Riemann Hypothesis}, Clay Mathematics Institute, 2000.
\bibitem{bgConnes99} A.~Connes, \emph{Trace formula in noncommutative geometry and the zeros of the Riemann zeta function}, Selecta Math.\ (N.S.) \textbf{5} (1999), 29--106.
\bibitem{bgdeBranges} L.~de Branges, \emph{Hilbert Spaces of Entire Functions}, Prentice-Hall, 1968.
\bibitem{bgMont73} H.~L.~Montgomery, \emph{The pair correlation of zeros of the zeta function}, Proc.\ Sympos.\ Pure Math.\ \textbf{24} (1973), 181--193.
\bibitem{bgBK} M.~V.~Berry and J.~P.~Keating, \emph{The Riemann zeros and eigenvalue asymptotics}, SIAM Rev.\ \textbf{41} (1999), 236--266.
\end{thebibliography}

\end{document}
